\documentclass[11pt,a4paper]{article}

\usepackage[T1]{fontenc}
\usepackage{lmodern}
\usepackage[margin=28mm]{geometry}
\usepackage{amsmath,amssymb,amsthm}
\usepackage{booktabs}
\usepackage{microtype}
\usepackage[hidelinks]{hyperref}

\newtheorem{theorem}{Theorem}
\newtheorem{lemma}[theorem]{Lemma}
\newtheorem{corollary}[theorem]{Corollary}
\newtheorem{proposition}[theorem]{Proposition}
\newtheorem{problem}[theorem]{Problem}
\theoremstyle{remark}
\newtheorem{remark}[theorem]{Remark}
\newtheorem{observation}[theorem]{Observation}
\newcommand{\Ostar}{O^{*}}
\newcommand{\E}{\mathbb{E}}
\newcommand{\W}{\mathcal{W}}
\newcommand{\nfree}{n_{\mathrm{free}}}

\clubpenalty=10000
\widowpenalty=10000
\displaywidowpenalty=10000
\emergencystretch=2em
\hypersetup{pdftitle={Max 2-SAT with Pruning Clauses: Pair Decisions, Sums, and a Measure}}

\title{Max 2-SAT with Pruning Clauses:\\ Pair Decisions, Sums, and a Measure}
\author{}
\date{6 October 2026, revised 7 October 2026\\[3pt]\small Working notes}

\begin{document}
\maketitle

\begin{abstract}
This note continues~\cite{N3}. It records the approach we now pursue, the
partial results, and the questions they raise; the revision of 7 October
incorporates results of the draft~\cite{N5}.
(i)~Pair decisions $u=v$ versus $u\ne v$ can be chained if every condition in
a $\ne$ child flips the whole class of variables tied together by earlier
decisions; flipping single variables loses optima, already for two clauses.
(ii)~Branching on variables of the active pruning clauses makes clauses cut
earlier but does not change how the saving decays with the degree.
(iii)~Adding up the active pruning clauses gives a sound cut rule in which
opposite literals of a shared variable cancel. On random formulas it saves
about $0.54n$ at every density we tried, up to $n=800$, while single-clause
checks fall to $0.19n$. Weighted sums cut exactly when the linear relaxation of
the active clauses is infeasible.
(iv)~For a measure-and-conquer analysis with weights on pruning clauses we
compute the optimal weights. Capped at one, the branching number depends on
the absolute allowance of the clause and tends to~$2$. Without the cap the
weights have a closed form, the clause length minus the logarithm of the number
of permitted completions, and the branching number depends only on the relative
allowance $\beta$ and tends to $2^{H(\beta)}$, which is tight.
(v)~The bound needs plain steps to be rare along one random walk, not on every
path.
(vi)~For formulas in which every variable has a linear gap relative to its
degree, we prove the parts that concern branching on near-dominated variables
and, for the rule without release, the resolution of one clause. What remains
is the accumulation over many clauses; in experiments plain steps cost at most a
constant factor, even on the circulant family.
(vii)~All these rules except pair decisions keep every strict local maximum, so
the number of strict local maxima bounds the size of their search trees from
below. Reaching $(2-\delta)^n$ at growing density therefore needs a rule that
moves several variables at once, such as class decisions.
\end{abstract}

\section{The approach}\label{sec:approach}

\paragraph{Goal.} As in~\cite{N3}: a deterministic algorithm for Max 2-SAT
with running time $\Ostar(c^n)$, $c<2$, in polynomial space, for formulas of
every clause density. The tools are DPLL with the domination
rule~\cite{N1,N2} and pruning clauses~\cite{N3}.

\paragraph{Plan.} We analyse the search with a measure in the spirit of
measure and conquer: the number of free variables minus a weight for each
counted pruning clause. A clause is counted only if it brings enough variables
that no other counted clause contains. The branching steps fall into three
kinds.
\begin{itemize}
 \item \emph{Case~1}: branch on a variable outside the counted clauses whose
 pruning clause will be counted. Variables close to domination give strong
 clauses.
 \item \emph{Case~2}: branch on a variable of a counted clause. In one child the
 clause gets shorter; in the other it gets tighter, and when its allowance is
 used up it propagates, fixing its remaining variables.
 \item \emph{Plain steps}: neither is possible.
\end{itemize}
The analysis is amortized over general scenarios rather than an exact case
analysis. It has a cold start: in a balanced formula the first pruning clauses
are weak, and strong ones appear only after many branchings. The plan is
therefore to analyse nodes with suitable structural properties and to show
that many nodes have them.

\paragraph{Three kinds of nodes.} The results below suggest a split.
\begin{enumerate}
 \item Many free variables are close to domination: the measure, with
 cases~1 and~2 (Sections~\ref{sec:measure}--\ref{sec:availability}).
 \item The formula is balanced, and the active clauses carry different
 literals on shared variables: sums of pruning clauses
 (Section~\ref{sec:sums}).
 \item Aligned structure, as in near-twins and the circulant family
 $C_{n,c}$ of~\cite{N3}: pair decisions (Section~\ref{sec:pairs}).
\end{enumerate}
Section~\ref{sec:barrier} sharpens this split: every rule except pair
decisions keeps all strict local maxima, so on formulas with many of them only
rules that move several variables at once can help.

\section{Setting}\label{sec:setting}

We use the notation of~\cite[Section~2]{N3} with $B=1$: a canonical formula
has at most one binary clause per pair. For a literal $\ell$ the deficit is
$p_\ell=d_2(\neg\ell)-d_1(\ell)+d_1(\neg\ell)$, the free degree of a variable
$v$ is $D_v=p_v+p_{\neg v}$, and $v$ is dominated when $\min\{p_v,p_{\neg v}\}\le
0$. We call $\min\{p_v,p_{\neg v}\}/D_v$ the \emph{ratio} of $v$; it is at most
$\frac12$.

When the search branches on $x$, let $\ell$ be the literal with the smaller
deficit $p=p_\ell\ge1$. The child with $\ell$ true is the \emph{preferred}
child. The other child, the \emph{pruned} child, receives the pruning clause
$P_x=\langle L_x\le p-1\rangle$, where $L_x$ consists of the literals $\neg m$
for the clauses $(\ell\lor m)$ and $m$ for the clauses $(\neg\ell\lor m)$.
By the flip identity~\cite[Lemma~18]{N3} an assignment that violates $P_x$ is
not worse after flipping $x$, and the search with domination, cuts and
propagation is exact for every branching rule~\cite[Proposition~20]{N3}.

\paragraph{Clause states.} At a node, a pruning clause is in \emph{state}
$(r,k)$ if $r$ of its literals are unassigned and at most $k$ of them may
still become true. Branching on $x$ creates the state $(D_x,p-1)$. When one of
its literals is assigned, the state moves to $(r-1,k)$ if the literal is false
and to $(r-1,k-1)$ if it is true. The clause is \emph{vacuous} when $k\ge r$,
it \emph{propagates} when $k=0$, setting all its remaining literals false, and
it is violated when $k<0$, in which case the node is cut.

\section{Pair decisions}\label{sec:pairs}

\subsection{One decision}

For a complete assignment $\alpha$ and a set $C$ of variables, let $\alpha^C$
be $\alpha$ with all variables of $C$ flipped and
$G_C(\alpha)=\mathrm{score}(\alpha^C)-\mathrm{score}(\alpha)$; write $G_u$ for
$G_{\{u\}}$. A \emph{pair decision} on unassigned variables $u,v$ has the
children $u=v$ and $u\ne v$. The first is explored in full. In the second we
keep only the assignments with $G_u(\alpha)<0$ and $G_v(\alpha)<0$; this is
sound, because if $G_u(\alpha)\ge0$, then $\alpha^{\{u\}}$ has $u=v$ and is at
least as good.

Let $\alpha$ have $u=1$ and $v=0$; the other case is symmetric. By the flip
identity, $G_u(\alpha)<0$ is the pruning clause of $u$ with respect to $\neg u$,
and $G_v(\alpha)<0$ is the pruning clause of $v$ with respect to $v$. Call a
common neighbour $w$ of $u$ and $v$ \emph{shared} if its clause with $u$ and
its clause with $v$ have the same signs, and let $\rho_u$, $\rho_v$ be the
numbers of the remaining clauses of $u$ and of $v$; a clause on $\{u,v\}$
counts for both. On a shared $w$ the two pruning clauses contain opposite
literals of $w$.

\begin{proposition}[Window]\label{prop:window}
Write $G_u(\alpha)=\sigma+a$ and $G_v(\alpha)=-\sigma+b$, where $\sigma$ is
the contribution of the clauses with shared neighbours and $a,b$ are the
contributions of the remaining clauses, so that $|a|\le\rho_u$ and
$|b|\le\rho_v$. The child $u\ne v$ keeps $\alpha$ only if
$b+1\le\sigma\le-a-1$. In particular:
\begin{enumerate}
 \item $a+b\le-2$, a condition on the non-shared clauses alone; it is the sum
 of the two pruning clauses, in which the shared part cancels;
 \item $\sigma$ takes at most $\rho_u+\rho_v-1$ values;
 \item if $\rho_u+\rho_v\le1$, the child is empty and $u,v$ can be merged.
\end{enumerate}
\end{proposition}

\begin{proof}
For a shared $w$, the clause of $w$ with $v$ at $v=0$ is the same function of
$w$ as the clause of $w$ with $u$ at $u=0$, and likewise at the value~$1$.
So flipping $v$ from $0$ to $1$ changes it by minus the amount by which
flipping $u$ from $1$ to $0$ changes the clause of $w$ with $u$. This gives the
form of $G_u$ and $G_v$; the rest is arithmetic.
\end{proof}

\subsection{Chaining decisions}

Along a path, decided pairs tie variables into classes. Write $[u]$ for the
class of $u$; the decisions fix every relation inside it, and flipping a whole
class preserves these relations.

\begin{proposition}[Chained pair decisions]\label{prop:chain}
Run the search of~\cite{N3} on classes: after a pair decision on $u,v$ the two
classes are merged in both children, with $u=v$ or $u=\neg v$; value
branchings, domination and pruning clauses apply to whole classes; and in the
child $u\ne v$ only the assignments with $G_{[u]}(\alpha)<0$ and
$G_{[v]}(\alpha)<0$ are kept, for the classes at the time of the decision.
Then the search is exact.
\end{proposition}

\begin{proof}
As in~\cite[Proposition~20]{N3}, order the complete assignments by the first
decision at which their paths part, preferring the child $u=v$ of a pair
decision and the preferred child of a value branching, and let $\alpha^*$ be
the largest optimal assignment. Suppose $\alpha^*$ is cut in the child $u\ne
v$ because $G_{[u]}(\alpha^*)\ge0$. Flipping all of $[u]$ changes no relation
inside $[u]$, and no earlier decision relates $[u]$ to another class, since
such a decision would have merged them; also no variable of $[u]$ is assigned
above the decision. So $\alpha^{*[u]}$ follows the path of $\alpha^*$ to the
decision, takes the child $u=v$ there, and is optimal, a contradiction. The
case $[v]$ and the other kinds of cuts are handled as in~\cite{N3}.
\end{proof}

Equivalently, a pair decision removes one free variable in each child, as a
value branching does, and the formula over merged variables turns a clause
inside a class into a unit clause or a constant.

\paragraph{Single flips are not enough.} Take $(x_0\lor x_1)$,
$(\neg x_1\lor\neg x_2)$ and decide $x_0?x_1$, then $x_1?x_2$. The optimum~$2$
is reached by $010$, $100$, $101$ and $110$ (values of $x_0x_1x_2$). The
assignment $110$ lies in the child $x_1\ne x_2$ below $x_0=x_1$. Flipping $x_1$
alone gives $100$, also optimal, so $110$ would be cut on the grounds that
$100$ lies in the child $x_1=x_2$. But $100$ has $x_0\ne x_1$ and lies below the
other child of the first decision, where flipping $x_1$ gives back $110$ and
cuts $100$. The two cuts point at each other. Also $010$ and $101$ are cut in
this way, and only $000$ and $111$, of score $1$, survive. Flipping the class
$\{x_0,x_1\}$ instead turns $110$ into $000$, a strict loss, and $110$ survives.

On chains $x_0?x_1$, $x_1?x_2,\dots$ in the circulant formulas $C_{9,2}$,
$C_{11,3}$, $C_{13,3}$, the single-flip version loses the optimum (it finds
$18$, $33$, $39$ instead of $28$, $43$, $60$), while the class version keeps
it, and only $10$, $10$ and $14$ complete assignments pass all its conditions.

\begin{remark}\label{rem:pairs-cnc}
In $C_{n,c}$ consecutive variables are near-twins: their environments differ
in two neighbours. That so few complete assignments pass suggests that pair
decisions attack the family that defeats every single-variable
rule~\cite[Theorem~13]{N3}. These are complete assignments, not tree nodes,
however: the class conditions involve the whole boundary of a class, and when
they can be checked during the search is open.
\end{remark}

\section{Branching strategies and timing}\label{sec:strategies}

\paragraph{Measuring trees.} For $n\ge50$ we estimate tree sizes by
population splitting~\cite{AV94}, a variant of Knuth's estimator~\cite{Knu75}:
keep $M$ random live nodes of each branching depth, expand all their children,
estimate the number of live nodes of the next depth, and resample. The
estimate is unbiased. On formulas with $n=30$ and $36$ it agrees with the exact
count within $0.3$ bits, already for $M=8$. We report the \emph{saving}
$n-\log_2(\text{tree size})$ as a fraction of $n$, averaged over two random
formulas, for random balanced formulas: each pair carries a clause with
probability $d/(n-1)$, with uniformly random signs.

\paragraph{Branching inside clauses.} The rule \emph{inside clauses} branches
on an unassigned variable of the active pruning clause that is most likely to
be violated, i.e.\ whose state $(r,k)$ maximizes $\Pr[\mathrm{Bin}(r,\frac12)
\ge k+1]$. Table~\ref{tab:strategies} compares it with a random static order.

\begin{table}[ht]
\centering
\begin{tabular}{lccccc}
\toprule
average degree $d$ & 16 & 32 & 64 & 128 & 200\\
\midrule
random order & 0.64 & 0.51 & 0.37 & 0.31 & 0.26\\
inside clauses & 0.80 & 0.64 & 0.48 & 0.39 & 0.30\\
random order and the sum rule (Section~\ref{sec:sums}) & 0.69 & 0.64 & 0.62
 & 0.54 & 0.54\\
\bottomrule
\end{tabular}
\caption{Saving as a fraction of $n$, for $n=400$.}\label{tab:strategies}
\end{table}

The strategy helps at every density, by about a quarter, but the saving still
decreases with $d$. For the random order it is close to
$1.55\,n\sqrt{(\ln d)/d}$, the timing estimate of~\cite[Remark~29]{N3}.

\begin{remark}[The timing in clause terms]\label{rem:timing}
A balanced variable with $D$ free neighbours and imbalance $M$, typically about
$\sqrt d$, gets a clause allowing about $(D-M)/2$ of its $D$ literals to be
true. Under random values about half of the assigned literals are true, so the
clause can cut only once about $D-M$ of its literals are assigned. Branching
inside a clause brings it there faster, but each such branching opens a new
clause whose support is almost entirely new: in a random formula two
neighbourhoods share about $d^2/n$ variables. So few clauses are near the end
of their supports at the same time. Random formulas have no near-twins, so
assuming that every two variables have different neighbourhoods does not remove
this effect.
\end{remark}

The \emph{coverage strategy}, which branches where the new clause brings the
most new variables and otherwise on clause variables adjacent to uncovered
ones, has not been tested yet.

\section{Sums of pruning clauses}\label{sec:sums}

Let $I$ be the set of active pruning clauses at a node that are not vacuous;
vacuous clauses can never be violated and would only weaken a sum. For an
unassigned variable $w$, let $m^1_w$ and $m^0_w$ be the numbers of clauses of
$I$ in which $w=1$, respectively $w=0$, makes a literal true, and let $k_i$ be
the allowance of clause $i$.

\begin{proposition}[Sum rule]\label{prop:sum}
If
\[
 \sum_{w\text{ unassigned}}\min\{m^0_w,m^1_w\}>\sum_{i\in I}k_i ,
\]
then no completion of the node satisfies all clauses of $I$, and the node can
be cut.
\end{proposition}

\begin{proof}
A completion that satisfies the clauses makes at most $\sum_ik_i$ further
literals of them true. Whatever value $w$ gets, it makes at least
$\min\{m^0_w,m^1_w\}$ of their literals true.
\end{proof}

The sum is implied by the clauses, so the search stays exact; we checked the
optimum against brute force on $150$ formulas with $n\le12$, in $42$ of which
the rule changed the tree. A variable shared with opposite literals contributes
in advance, whatever its value; one shared with the same literal contributes
nothing. This is the algorithmic side of~\cite[Lemma~27]{N3}: shared opposite
literals give a smaller joint survival probability than independent clauses.
Proposition~\ref{prop:window}(1) is the case of two clauses.

\paragraph{Example.} $\langle\{w_1,\dots,w_7\}\le3\rangle$ and
$\langle\{\neg w_1,\dots,\neg w_7\}\le3\rangle$ say that at most three of the
$w_i$ are $1$ and at most three are $0$, which is impossible. Neither clause is
violated or propagates at the start. A search with single-clause checks and
propagation branches until three equal values appear and has $20$ leaves. In
the sum every $w_i$ contributes $\min\{1,1\}=1$, so $7>3+3$ and the node is
cut at once.

\begin{remark}[Why the sum decides early]\label{rem:sum-heur}
This is a heuristic count. Suppose a path carries $k$ active clauses, created at
variables of imbalance about $\sqrt d$, so that every bound lies about
$\sqrt d/2$ below the expected count of its clause. These shortfalls point the
same way and add up to order $k\sqrt d$. In a random formula an unassigned $w$
lies in about $m_w=kd/n$ of the clauses, with independent literal signs, so the
uncertain part $\sum_w|m^1_w-m^0_w|$ over the $u$ unassigned variables is of
order $u\sqrt{kd/n}$. The sum decides once this falls below the shortfall,
i.e.\ once $u\lesssim\sqrt{kn}$. Since $k$ grows in proportion to the number
of assigned variables, this happens while a constant fraction of the
variables is unassigned, and $d$ cancels.
\end{remark}

\begin{table}[ht]
\centering
\begin{tabular}{lcccccc}
\toprule
$n$ & 50 & 100 & 200 & 400 & 600 & 800\\
\midrule
single-clause checks & 0.48 & 0.39 & 0.29 & 0.26 & 0.22 & 0.19\\
with the sum rule & 0.53 & 0.44 & 0.51 & 0.54 & 0.54 & 0.58\\
\bottomrule
\end{tabular}
\caption{Saving as a fraction of $n$ on dense random formulas ($d\approx
n/2$), random order. For $n=600$ and $800$: one formula and $M=40$.}
\label{tab:dense}
\end{table}

With single-clause checks the saving grows roughly like $2\sqrt{n\ln n}$; with the sum
rule it stays near $n/2$ (Table~\ref{tab:dense}), and across the degrees of
Table~\ref{tab:strategies} it flattens at about $0.54n$. On the aligned
formulas $C_{n,c}$, $15\le n\le21$, which have no opposite literals, the sum
removes at most $5\%$ of the nodes.

Any nonnegative combination of the active clauses is also sound, and the
family of all of them has an exact description~\cite{N5}.

\begin{proposition}[Sums are linear-relaxation cuts]\label{prop:lp}
Some nonnegative combination of the active clauses $L_i\le k_i$ cuts a node,
in the sense of Proposition~\ref{prop:sum}, if and only if the linear
relaxation $\{x\in[0,1]^r:L_i(x)\le k_i\text{ for all }i\}$ is empty.
\end{proposition}

\begin{proof}
With weights $a_i\ge0$ the rule cuts iff the minimum of $\sum_ia_iL_i$ over
$\{0,1\}^r$ exceeds $\sum_ia_ik_i$. The function is affine, so this is also its
minimum over $[0,1]^r$, and a point of the relaxation rules the cut out.
Conversely, if the relaxation is empty, the image of the box under
$x\mapsto(L_i(x))_i$, a compact convex set, misses the closed convex set
$\{y:y_i\le k_i\}$. A separating hyperplane must have nonnegative coefficients,
since the second set is unbounded below in every coordinate, and these
coefficients give a cut.
\end{proof}

So sums see exactly what the linear relaxation sees, although the separating
weights may be large. Contradictions of parity type are invisible to them:
$x+y=y+z=z+x=1$, written as six cardinality clauses, has the fractional solution
$(\frac12,\frac12,\frac12)$ but no Boolean one~\cite{N5}. Counting the
assignments that satisfy a weighted sum, instead of only testing it, gives a
graded certificate that satisfies a branching inequality, together with a
sufficient condition for a running-time bound in terms of the progress of this
count along a random walk~\cite{N5}. In these terms Remark~\ref{rem:sum-heur}
says, heuristically, that on random formulas the relaxation of the active
clauses becomes infeasible while about half of the variables are free.

\section{A measure for pruning clauses}\label{sec:measure}

\subsection{Counting, measure and rule}

A clause is counted only on its new variables: when it is created, delete the
literals of variables that lie in the support of a counted clause. Deleting
literals keeps a pruning clause valid, since it can only lower the count, so
the counted clauses are valid and have pairwise disjoint supports. The search
may keep the full clauses for cuts; the measure uses only the counted ones. Let
\[
 \mu=\nfree-\sum_{\text{counted }C}\varepsilon(r_C,k_C),
\]
where $(r_C,k_C)$ is the state of $C$ and $\varepsilon\ge0$ is a weight to be
chosen. The rule:
\begin{enumerate}
 \item if a counted clause is active, branch on one of its variables
 (case~2);
 \item otherwise, if a free variable has a \emph{feasible} new clause (defined
 below), branch on it (case~1);
 \item otherwise branch on any variable (a plain step).
\end{enumerate}
A counted clause in a state of weight $0$ is \emph{released}: it no longer
counts.

\subsection{Steps and weights}

A step is \emph{admissible} for $\lambda\in(1,2)$ if its children have measure
drops $d_0,d_1$ with $\lambda^{-d_0}+\lambda^{-d_1}\le1$.
\begin{itemize}
 \item Case~1 on $x$ with new degree $D$ and deficit $p$: drops $1$ and
 $1+\varepsilon(D,p-1)$, or $1$ and $1+D$ if $p=1$. This is admissible iff
 $\varepsilon(D,p-1)\ge c(\lambda):=\log_\lambda\frac1{\lambda-1}$; we call
 the new clause feasible in that case.
 \item Case~2 in state $(r,k)$: the false child has drop
 $1+V_f-\varepsilon(r,k)$ with $V_f=\varepsilon(r-1,k)$, or $V_f=0$ if
 $(r-1,k)$ is vacuous or released; the true child has drop
 $1+V_t-\varepsilon(r,k)$ with $V_t=\varepsilon(r-1,k-1)$, and $V_t=r-1$ if
 $k=1$, when the other $r-1$ literals are set false.
 \item A plain step has drops of at least $1$.
\end{itemize}
Further variables fixed in a child only add to its drop, as long as no forced
transition lowers a weight by more than one.

\begin{lemma}[Optimal weights]\label{lem:weights}
For $\lambda\in(1,2)$ and a cap $\kappa\in(0,\infty]$ define, for $r=2,3,\dots$
and $1\le k<r$,
\[
 \varepsilon(r,k)=\max\Bigl\{0,\ \min\bigl\{\kappa,\
 1-\log_\lambda\bigl(\lambda^{-V_f}+\lambda^{-V_t}\bigr)\bigr\}\Bigr\},
\]
with $V_f,V_t$ as above. These are the largest weights, bounded by $\kappa$,
for which every case-2 step in a state of positive weight is admissible.
Moreover $\varepsilon(r,k)\le1+\min\{V_f,V_t\}$, so no transition lowers a
weight by more than one, forced assignments never increase $\mu$, and
$\varepsilon(r,k)\le r-1$, so $\mu\ge0$.
\end{lemma}

\begin{proof}
The step in state $(r,k)$ is admissible iff
$\lambda^{\varepsilon(r,k)}(\lambda^{-1-V_f}+\lambda^{-1-V_t})\le1$, which is the
bound on $\varepsilon(r,k)$; it depends only on shorter states, and larger
children allow a larger parent. Since $\lambda^{-V_f}+\lambda^{-V_t}\ge
\lambda^{-\min\{V_f,V_t\}}$, we get $\varepsilon(r,k)\le1+\min\{V_f,V_t\}$; a
forced literal then changes $\mu$ by $-1$ plus a weight decrease of at most
one. Induction on $r$ with $V_f\le r-2$ gives $\varepsilon(r,k)\le r-1$.
\end{proof}

\subsection{With and without the cap}

\begin{proposition}[Cap one]\label{prop:cap}
With $\kappa=1$, for large $r$ the weights approach $1,1-c,1-2c,\dots$ for
$k=1,2,3,\dots$, where $c=c(\lambda)$. So case~1 with allowance $k_0=p-1$ is
admissible for large $D$ iff $\lambda^{1/k_0}(\lambda-1)\ge1$.
\end{proposition}

\begin{proof}
For large $r$, $V_t=r-1\to\infty$ in state $(r,1)$, so its weight is the cap
$1$. For $k\ge2$ a stationary weight $\varepsilon_k$ must satisfy
$\lambda^{1-\varepsilon_k}\ge\lambda^{-\varepsilon_k}+\lambda^{-\varepsilon_{k-1}}$,
i.e.\ $\varepsilon_k\le\varepsilon_{k-1}-c$. Case~1 needs
$1-(k_0-1)c\ge c$, i.e.\ $k_0c\le1$.
\end{proof}

\begin{table}[ht]
\centering
\begin{tabular}{lcccccc}
\toprule
allowance $k_0=p-1$ & 1 & 2 & 3 & 4 & 8 & 16\\
\midrule
$\lambda$ with cap one & 1.618 & 1.756 & 1.820 & 1.857 & 1.922 & 1.959\\
\bottomrule
\end{tabular}
\caption{The cap $\varepsilon<1$: the branching number depends only on $k_0$.}
\label{tab:cap}
\end{table}

So with the cap $\lambda$ tends to $2$, roughly as $2-(\ln2)/k_0$: the $k_0$
units of allowance share one unit of weight. Only variables a bounded number of units away from
domination help, and closeness to domination in the relative sense,
$p\le\beta D$, gives nothing for large degrees.

Without the cap the weights have a closed form~\cite{N5}. Let
$V(r,k)=\sum_{j=0}^{\min\{r,k\}}\binom rj$ be the number of completions that a
clause in state $(r,k)$ permits.

\begin{theorem}[No cap]\label{thm:nocap}
The weights $e_\lambda(r,k)=r-\log_\lambda V(r,k)$ satisfy the recursion of
Lemma~\ref{lem:weights} with equality and without the clipping at zero, if a
vacuous state keeps its value $e_\lambda(r,k)=r(1-\log_\lambda2)<0$ instead of
$0$. With these weights every case-2 step is tight, no transition lowers a
weight by more than one, and creating a clause in state $(D,k_0)$ is admissible
if and only if
\[
 V(D,k_0)\le(\lambda-1)\lambda^D .
\]
Hence, for $k_0\le\beta D$ with $\beta<\frac12$ and any $\lambda>2^{H(\beta)}$,
every creation with $D\ge\ln\frac1{\lambda-1}\big/\bigl(\ln\lambda-H(\beta)\ln2
\bigr)$ is admissible, and the smallest admissible $\lambda$ for
$k_0=\lfloor\beta D\rfloor$ tends to $2^{H(\beta)}$ as $D\to\infty$.
\end{theorem}

\begin{proof}
By Pascal's rule $V(r,k)=V(r-1,k)+V(r-1,k-1)$, so
$\lambda^{-e_\lambda(r-1,k)}+\lambda^{-e_\lambda(r-1,k-1)}=\lambda^{1-e_\lambda(r,k)}$,
which is the recursion with equality; for $k=1$ the true child has
$V(r-1,0)=1$, i.e.\ weight $r-1$, as for propagation. A transition changes the
weight by $-1+\log_\lambda\frac{V(r,k)}{V(r-1,k')}\ge-1$, since a restriction
cannot increase $V$. Case~1 is admissible iff
$\lambda^{-1}+\lambda^{-1-D}V(D,k_0)\le1$, which is the stated condition. The
threshold follows from $V(D,k_0)\le2^{H(\beta)D}$, and the limit from
$V(D,k_0)\ge\binom D{k_0}\ge2^{H(k_0/D)D}/(D+1)$.
\end{proof}

The weights of Lemma~\ref{lem:weights} with $\kappa=\infty$ are these weights
with \emph{release}: where $e_\lambda$ would be negative the clause is released
and its weight is $0$. Clipping only raises weights, so it can only lower the
branching numbers, slightly at finite $D$ (Table~\ref{tab:nocap}). With release,
every creation with $k_0\le\beta D$ and $D\le400$ is already admissible at
$\lambda=2^{H(\beta)}$, for $\beta\in\{\frac18,\frac14,0.3,\frac38,0.45\}$.

\begin{table}[ht]
\centering
\begin{tabular}{llcccccc}
\toprule
$\beta$ & & $D=32$ & $64$ & $128$ & $256$ & $512$ & $2^{H(\beta)}$\\
\midrule
$1/8$ & with release & 1.411 & 1.419 & 1.429 & 1.439 & 1.446 & 1.458\\
 & without & 1.431 & 1.437 & 1.443 & 1.448 & 1.452 & \\
$1/4$ & with release & 1.668 & 1.693 & 1.713 & 1.728 & 1.739 & 1.755\\
 & without & 1.695 & 1.715 & 1.730 & 1.740 & 1.746 & \\
$3/8$ & with release & 1.849 & 1.877 & 1.898 & 1.913 & 1.923 & 1.938\\
 & without & 1.873 & 1.897 & 1.912 & 1.923 & 1.929 & \\
\bottomrule
\end{tabular}
\caption{Without the cap: the branching number for creating a clause with
$k_0=\beta D$, with release (Lemma~\ref{lem:weights}) and without
(Theorem~\ref{thm:nocap}).}\label{tab:nocap}
\end{table}

For example, with release at $\lambda=1.643$ the weights of the states $(16,k)$ are
$10.34$, $6.22$, $3.10$, $0.89$ for $k=1,\dots,4$, those of $(12,k)$ are
$6.89$, $3.38$, $1.00$, and those of $(8,k)$ are $3.66$, $1.03$. At fixed
allowance a longer clause weighs more; at fixed margin $r-2k$ a shorter one
does.

\begin{remark}[Tightness]\label{rem:tight}
If every variable has ratio $\beta$, the bound $2^{H(\beta)n}$ is that
of~\cite[Theorem~11]{N3} with $\beta=(1-\varepsilon)/2$, and it is attained on
$C_{n,c}$ by every branching rule with domination and pruning
clauses~\cite[Theorem~13, Observation~22]{N3}. So the step analysis is optimal
for this class.
\end{remark}

\begin{remark}[Why one might cap, and why not to]\label{rem:onebit}
In base $2$ a disjoint block saves exactly $\log_2\frac2{1+\theta}<1$
bit~\cite[Proposition~25]{N3}, since the preferred child carries no clause: a
clause is worth less than one bit to the whole tree. This suggests $\varepsilon<1$,
so that the variable assigned when a clause disappears pays for it. In an
analysis by branching numbers, however, the weight must also pay for the steps
inside the clause before it propagates, about one per unit of allowance. Without
the cap the weight of a tight clause is close to the number of variables it will
force: $(16,1)$ weighs $10.34$ and forces $15$ when it propagates. What the
payoff argument needs is Lemma~\ref{lem:weights}: no transition lowers a weight
by more than one.
\end{remark}

\begin{remark}[What nearly disjoint clauses can certify]\label{rem:disjointcap}
Heuristically, for a random formula of degree $d$: a clause created when its
variable has $D$ free neighbours is worth $O(1+d/D)$, since the imbalance is
$O(\sqrt d)$; at most $\nfree/D$ disjoint counted clauses coexist; and
$D\approx d\,\nfree/n$. While $\nfree$ falls from $2N$ to $N$, the counted
clauses therefore certify $O\bigl((n/d)(n/N)\bigr)$, which pays for all $N$
variables only once $N\lesssim n/\sqrt d$. Summing over the stages gives
$O(n/\sqrt d)$ in total, the timing barrier again. Counting only nearly disjoint
clauses cannot give $c^n$ for balanced formulas of growing degree.
\end{remark}

\section{Averaging over paths}\label{sec:averaging}

\begin{lemma}[Averaging]\label{lem:avg}
Let $\lambda\in(1,2]$, let the measure be $n$ at the root and nonnegative
everywhere, and let every internal node be either admissible or plain, as in
Section~\ref{sec:measure}. Let $\W$ be the random walk from the root that at an
admissible node moves to child $i$ with probability $q_i\ge\lambda^{-d_i}$,
where $q_0+q_1\le1$ and the walk stops with the remaining probability, and at a
plain node moves to each child with probability $\frac12$. Let $S$ be the
number of plain nodes on $\W$. Then
\[
 \#\text{leaves}\le\lambda^n\,\E\bigl[(2/\lambda)^S\bigr].
\]
\end{lemma}

\begin{proof}
For a leaf $\ell$ let $s(\ell)$ be the number of plain nodes on its path and
$\Pi(\ell)$ the probability that $\W$ ends at $\ell$. The drops along the path
add up to at most $n$, those at plain nodes to at least $s(\ell)$, so the drops
at admissible nodes add up to at most $n-s(\ell)$ and
$\Pi(\ell)\ge\lambda^{-(n-s(\ell))}2^{-s(\ell)}$. Hence
$\#\text{leaves}=\sum_\ell\Pi(\ell)/\Pi(\ell)\le\lambda^n\sum_\ell\Pi(\ell)
(2/\lambda)^{s(\ell)}\le\lambda^n\E[(2/\lambda)^S]$.
\end{proof}

\begin{corollary}\label{cor:avg}
\begin{enumerate}
 \item If every path has at most $s$ plain nodes, there are at most
 $\lambda^{n-s}2^s$ leaves. With the weights of Lemma~\ref{lem:weights} and
 $\lambda=2^{H(\beta)}$, this is the bound for the rule of
 Section~\ref{sec:measure} when every node offers case~1 with $k_0\le\beta D$
 or case~2, for $D\le400$; for larger $D$ take $\lambda$ slightly above
 $2^{H(\beta)}$ (Theorem~\ref{thm:nocap}).
 \item If $\Pr[S>s]\le\rho$, there are at most $\lambda^{n-s}2^s+\rho\,2^n$
 leaves, since a walk has at most $n$ plain nodes.
 \item If $S-\bar s$ is sub-Gaussian with variance proxy $\sigma^2$, there are
 at most $\lambda^n(2/\lambda)^{\bar s}e^{\theta^2\sigma^2/2}$ leaves, where
 $\theta=\ln(2/\lambda)$.
\end{enumerate}
\end{corollary}

Item~2 has the shape of the bound of~\cite{N2}: most branches are good, and the
bad ones are exponentially rare. Two points matter.
\begin{itemize}
 \item The probability is over the tilted walk $\W$, not over uniform
 assignments. At a case-1 step $\W$ takes the preferred child with probability
 $1/\lambda$, and it may give the pruned child $1-1/\lambda$, which is at least
 $\lambda^{-1-\varepsilon}$ exactly when the step is admissible.
 \item The failure probability must be exponentially small, because bad paths
 together can cost $2^n$. In item~3 the price of fluctuations is second order:
 for $\lambda=1.755$ and $\sigma^2=n/4$ the factor is about $2^{0.003n}$.
\end{itemize}
Variables fixed during plain steps can be credited at rate $\lambda$: if
$\pi(\ell)\ge s(\ell)$ is the total drop at the plain nodes of the path, the
proof gives $\#\text{leaves}\le\E[\lambda^{n-\pi}2^S]$. This is how
domination after plain steps, the mechanism of~\cite[Theorem~11]{N3}, enters.

\section{Availability in formulas with linear gaps}\label{sec:availability}

Assume every variable starts with ratio at most $\beta_0<\frac12$, fix a
threshold $\beta'\in(\beta_0,\frac12)$, let $\lambda=2^{H(\beta')}$, and let case~1
take the free variable of smallest ratio among those with $p-1\le\beta'D$ and
a feasible clause. Since active clauses are finished first, a plain step
happens only when no free variable passes the threshold. We want
$\E[(2/\lambda)^S]=2^{o(n)}$.

Each assignment of a neighbour $y$ of $v$ lowers exactly one of the two
deficits of $v$~\cite[Lemma~2]{N3}. Call it \emph{helpful} for $v$ if it
lowers the smaller one.

\begin{lemma}[Window]\label{lem:window}
Let $v$ start with degree $D_0$ and smaller deficit $p_0\le\beta_0D_0$. If at
least a $\beta_0$-fraction of its assigned neighbours were helpful, then the
ratio of $v$ is at most $\beta_0$.
\end{lemma}

\begin{proof}
With $j$ assigned neighbours, $h\ge\beta_0j$ of them helpful, the smaller
deficit is at most $p_0-h\le\beta_0(D_0-j)$ and the degree is $D_0-j$.
\end{proof}

Conversely, assignments that are not helpful lower the other deficit, and once
enough of them have occurred, the other literal of $v$ passes the threshold.
For $p_0=\beta_0D_0$ and no helpful assignments this happens after a fraction
$(1-\beta_0-\beta')/(1-\beta')$ of the neighbours, $64\%$ for $\beta_0=\frac14$
and $\beta'=0.3$. So $v$ fails the threshold only inside a bounded window, and
a plain step needs every free variable inside its window at the same time.

\begin{lemma}[Case-1 steps]\label{lem:coins}
Let $\W$ take the preferred child of a case-1 step with probability $1/\lambda$
and the pruned child with probability $1-1/\lambda$. Then every neighbour of
the branched variable is helped with probability at least $1-1/\lambda$, and
$1-2^{-H(\beta)}\ge\beta$ for $0<\beta\le\frac12$, with equality only at
$\beta=\frac12$.
\end{lemma}

\begin{proof}
Exactly one value of the branched variable helps a given neighbour, and both
values have probability at least $1-1/\lambda$. The inequality is
$H(\beta)\ge\log_2\frac1{1-\beta}$, which after expanding $H$ reads
$\beta\log_2\frac{1-\beta}\beta\ge0$.
\end{proof}

With $\lambda=2^{H(\beta')}$ each case-1 step helps every neighbour with
probability at least $\beta'>\beta_0$. If all assignments were case-1 steps,
the Azuma--Hoeffding inequality would bound the probability that a variable
with $m$ assigned neighbours has a helpful fraction below $\beta_0$ by
$e^{-2(\beta'-\beta_0)^2m}$, and Lemma~\ref{lem:window} would keep it at the
threshold.

\paragraph{The obstacle.} In a pruned child, $\W$ then resolves the clause of
the branched variable $x$, which allows at most a $\beta'$-fraction of the
neighbours of $x$ to take their bad values. These assignments follow the
clause, not a coin. In aligned formulas, where after renaming every clause is
$(x\lor y)$ or $(\neg x\lor\neg y)$, the bad value of $y$ for the clause of
$x$ is its minority value, which is the helpful value for every other neighbour
of $y$. So the other variables see a helpful fraction of about the ratio of
$x$, which is below their own if $x$ has the smallest ratio. The per-variable
argument does not cover this, and on $C_{n,c}$ plain steps do occur.

\paragraph{Measurements.} On random formulas with linear gaps, with biased
signs or aligned, $n\in\{200,400\}$, $d$ from $20$ to $200$, $\beta_0=\frac14$
and $\beta'=0.26$, no walk had a plain step. Even with $\beta'=\beta_0$ on dense
aligned formulas, a walk had at most $16$ plain steps and
$\log_2\E[(2/\lambda)^S]=0.10$. Table~\ref{tab:cnc} shows the extreme case.

\begin{table}[ht]
\centering
\begin{tabular}{lcccc}
\toprule
 & \multicolumn{2}{c}{$\beta'-\beta_0=0.02$} &
 \multicolumn{2}{c}{$\beta'-\beta_0=0.05$}\\
\cmidrule(lr){2-3}\cmidrule(lr){4-5}
$n$ & walks with a plain step & $\log_2\E[(2/\lambda)^S]$ &
walks with a plain step & $\log_2\E[(2/\lambda)^S]$\\
\midrule
101 & 20\% & 0.49 & 9\% & 0.29\\
201 & 19\% & 1.32 & 12\% & 0.17\\
401 & 23\% & 2.01 & 17\% & 0.55\\
801 & 26\% & 2.73 & 11\% & 0.63\\
\bottomrule
\end{tabular}
\caption{Plain steps along $\W$ on $C_{n,c}$ with $\varepsilon=0.5$
($\beta_0\approx\frac14$), $150$ walks ($80$ for $n=801$).}
\label{tab:cnc}
\end{table}

So the bound is $\lambda^n$ up to a factor below $8$, that is
$2^{H(\beta')n}$. For gaps that are linear in $n$ this exponent is already
known for every order~\cite[Theorems~11 and~24]{N3}; what is new here is that
the gaps are relative to each variable's own degree, so the formulas may have
any density, and that the rule is adaptive.

\paragraph{Why, heuristically.} In random formulas the variables reach their
windows at different times, and the rule takes the best free variable, so one
outside its window is always available. In $C_{n,c}$ all variables move
together, but $\W$ follows typical surviving assignments. For a uniformly
random assignment satisfying a clause, the remaining allowance divided by the
remaining length is a martingale, so the bad fraction stays near the ratio of
the clause, the minority fraction stays near $\beta_0$, and only fluctuations
reach the window.

\begin{problem}[The missing lemma]\label{prob:phase}
Show that while $\W$ resolves a counted clause created in state $(D,k_0)$, the
fraction of its literals that become true is at least $k_0/D-\eta$, except with
probability $e^{-\Omega(\eta^2D)}$.
\end{problem}

For the rule with release this is open. Without release it holds, in a
stronger form.

\begin{proposition}[One clause without release]\label{prop:phase}
Use the weights $e_\lambda$ of Theorem~\ref{thm:nocap}, keep every counted
clause until its support is assigned, and postpone domination and new clauses
while it is active. Then inside the clause $\W$ makes all $V(D,k_0)$ permitted
completions equally likely, and for $0\le k_0<D/2$, $1\le t\le k_0$ and
$q=k_0/(D-k_0+1)$ the number $J$ of its literals that become true satisfies
\[
 \Pr[J\le k_0-t]\le q^t .
\]
\end{proposition}

\begin{proof}
With tight weights, $\W$ moves to a child with probability
$\lambda^{-\text{drop}}=V(\text{child})/V(\text{parent})$, and these factors
telescope to $1/V(D,k_0)$ along every permitted completion. For $j\le k_0$,
$\binom D{j-1}=\frac j{D-j+1}\binom Dj\le q\binom Dj$, hence
$\binom Di\le q^t\binom D{i+t}$ for $i\le k_0-t$; summing over $i$ and dividing
by $V(D,k_0)$ gives the bound~\cite{N5}.
\end{proof}

So Problem~\ref{prob:phase} holds for this rule, with $e^{-\Omega(\eta D)}$ in
place of $e^{-\Omega(\eta^2D)}$: $J$ ends within $O(1)$ of the allowance. Given
$J$, the true literals form a uniformly random subset of the support. Hence a
variable with $s$ neighbours in the support, for each of which one value is
designated as helpful when the clause is created, sees a helpful fraction of at
least $k_0/D-\eta$ among them, whatever the designations, except with
probability $q^{\lceil\eta D/2\rceil}+2e^{-\eta^2s/8}$~\cite{N5}. The price of
giving up release is a slightly larger branching number (Table~\ref{tab:nocap}).

This suggests how to finish the argument. By Lemma~\ref{lem:window} a variable
stays at the threshold while at least a $\beta_0$-fraction of its assigned
neighbours helped it; case-1 steps help with probability at least $\beta'$
(Lemma~\ref{lem:coins}), and a clause passes on a helpful fraction of at least
its ratio minus $\eta$. So in case~1 we should take the largest ratio below
$\beta'$, not the smallest: when this ratio is close to $\beta'$, every helpful
fraction exceeds $\beta_0$, and the creation is still feasible. What remains
open is the accumulation of these bounds over many adaptively chosen clauses,
in particular for small overlaps $s$, and the assignments that domination forces
between clauses. We have not tested the largest-ratio choice; on random
formulas there are no plain steps to remove, and in $C_{n,c}$ all variables have
the same ratio.

\section{A barrier for single-flip rules}\label{sec:barrier}

Call a complete assignment a \emph{strict local maximum} if flipping any single
variable strictly lowers the score.

\begin{proposition}[Strict local maxima are never cut]\label{prop:barrier}
Consider a search whose fixings and cuts are domination, pruning clauses with
propagation, and sums or weighted sums of active pruning clauses; the full
phases and counting certificates of~\cite{N5} are included. For every branching
rule, every strict local maximum is a leaf of its search tree. Hence the tree
has at least as many leaves as the formula has strict local maxima.
\end{proposition}

\begin{proof}
Let $\alpha$ be a strict local maximum and follow its path. If domination fixes
a literal $\ell$, flipping $\alpha$ towards $\ell$ does not lower the score by
the flip identity, so $\alpha$ already has $\ell$ true. If $\alpha$ enters a
pruned child, flipping the branched variable strictly lowers the score, so by
the flip identity $\alpha$ satisfies the new pruning clause. Hence $\alpha$
satisfies every clause on its path, agrees with every propagation they cause,
and satisfies every sum of them and every count that includes their common
solutions.
\end{proof}

\paragraph{Many strict local maxima.} Max-Cut on $K_m$, written with the clauses
$(x_i\lor x_j)$ and $(\neg x_i\lor\neg x_j)$ for every pair, has
$\binom m{m/2}$ strict local maxima for even $m$, the balanced cuts. Disjoint
copies with $d=m-1$ have $2^{\,n-O(n(\log d)/d)}$ of them, which is the order of
the saving of~\cite{GK14}. Disjoint copies are separated by splitting into
components, but Hastings shows that his upper bounds on the number of local
maxima are tight, up to polylogarithmic factors for dense formulas and up to a
factor of order $d$ for sparse ones~\cite{Has16}. So single-flip reasoning,
however refined, cannot by itself give $(2-\delta)^n$ at growing density. In
particular the progress condition of~\cite{N5} fails on such formulas, since the
count at every depth is at least the number of strict local maxima.

\paragraph{What can cut them.} Pair decisions can: in a $\ne$ child they compare
an assignment with the flip of a whole class, a move of several variables. On
$K_8$, $K_{10}$ and $K_{12}$ the chained class conditions of
Section~\ref{sec:pairs} leave $36$, $94$ and $276$ complete assignments,
against $70$, $252$ and $924$ strict local maxima. For exact twins this is
visible directly: merged with $u=\neg v$, the two clauses of a twin pair with a
common neighbour $w$ add up to one plus the indicator of the literal of $w$,
independent of $u$. So the merged variable keeps only the clauses between $u$
and $v$, domination fixes it, and the two orientations of the pair are never
both explored. Hastings also constructs sparse formulas with many maxima under
moves of up to $k$ variables~\cite{Has16}, so moves of bounded size may meet a
similar barrier; class flips are not of bounded size.

\paragraph{A second obstacle.} The circulant formulas $C_{9,2}$, $C_{11,3}$ and
$C_{13,3}$ have no strict local maxima at all, yet on $C_{n,c}$ every rule
of~\cite{N3} costs $\binom{n-1}a$ leaves~\cite[Theorem~13, Observation~22]{N3}.
Their leaves are only one-sided local maxima: the preferred child is never
pruned. The class conditions reduce these three formulas to $10$, $10$ and $14$
complete assignments.

\paragraph{Consequence for the plan.} On formulas with few strict local maxima,
single-flip rules and their sums can in principle suffice, and the question is
when they cut (Sections~\ref{sec:strategies} and~\ref{sec:sums}). On formulas
with many, or where one-sided pruning is the bottleneck, a rule that moves
several variables at once is necessary, and among our tools only class
decisions qualify.

\section{Numerical checks}\label{sec:numerics}

The scripts are in this folder.
\begin{itemize}
 \item \texttt{explore\_pair\_chain.py} (seconds): the two-clause example, the
 chains in $C_{n,c}$ and in Max-Cut on $K_m$, and the numbers of strict local
 maxima (Sections~\ref{sec:pairs} and~\ref{sec:barrier}).
 \item \texttt{explore\_sum\_rule.py}: \texttt{check} (seconds) compares the
 optimum with brute force on $150$ formulas, with and without the sum rule and
 for both strategies, and the estimator with exact counts for $n=30,36$;
 \texttt{density} and \texttt{dense} (tens of minutes) give
 Tables~\ref{tab:strategies} and~\ref{tab:dense} up to $n=400$.
 \item \texttt{explore\_clause\_weights.py} (two seconds):
 Proposition~\ref{prop:cap}, the rows with release of Table~\ref{tab:nocap},
 the weight table, and the check that no transition lowers a weight by more
 than one. The rows without release solve $V(D,k_0)=(\lambda-1)\lambda^D$
 (Theorem~\ref{thm:nocap}).
 \item \texttt{explore\_availability.py} (ten seconds): the measurements of
 Section~\ref{sec:availability} for $n\le401$.
\end{itemize}
The entries for $n=600,800$ in Table~\ref{tab:dense} and for $n=801$ in
Table~\ref{tab:cnc} are single runs of the same code with larger $n$.

\section{Open questions and new directions}\label{sec:open}

\begin{enumerate}
 \item \textbf{Accumulation over clauses.} Combine
 Proposition~\ref{prop:phase} with Lemmas~\ref{lem:window} and~\ref{lem:coins}
 over many adaptively chosen clauses, including the assignments forced by
 domination between them, or settle Problem~\ref{prob:phase} for the rule with
 release. This would give $2^{H(\beta')n}$, for every $\beta'>\beta_0$ and every
 density, for formulas in which every variable has a gap linear in its degree,
 with an adaptive rule.
 \item \textbf{Partial gaps.} If only $\eta n$ variables have gaps, plain steps
 on the others are unavoidable. Does the refined form of Lemma~\ref{lem:avg},
 which credits domination during plain steps, recover the exponent
 $\max\{1-\eta,H((1-\varepsilon)/2)\}$ of~\cite[Theorem~11]{N3}?
 \item \textbf{Ending the cold start.} Find structural properties that produce
 near-dominated variables in balanced formulas. A candidate is coverage: a set
 $S$ whose variables have most of their neighbours outside $S$. After a random
 assignment of the rest, a variable of $S$ is dominated once its imbalance, a
 $\pm1$ walk over its outside neighbours, exceeds its degree inside $S$; by the
 Erd\H{o}s--Littlewood--Offord bound~\cite{Erd45} this fails with probability
 $O(b/\sqrt D)$ for inside degree $b$ and outside degree $D$. So inside degree
 up to $\varepsilon\sqrt D$ suffices, not only the bounded degree needed for
 independence~\cite{N5}, provided the fields of $S$ are jointly spread out;
 near-twins have identical fields, which is where that assumption is needed.
 This is the idea behind Hastings' bound~\cite{Has16}, and in pseudo-random
 formulas such sets have size $O(n/\sqrt d)$, so coverage alone gives
 degree-dependent savings.
 \item \textbf{Counting overlapping clauses.} By Remark~\ref{rem:disjointcap}
 the measure must count groups of overlapping clauses. The weighted-sum counts
 of~\cite{N5} provide such a joint weight, with a branching inequality and a
 running-time bound under a condition on their progress. The open part is to
 show linear progress before the last levels for an interesting class, starting
 with pseudo-random formulas.
 \item \textbf{The sum rule.} Prove the behaviour of Table~\ref{tab:dense} on
 random formulas (Remark~\ref{rem:sum-heur}); try the best weighted sums, which
 by Proposition~\ref{prop:lp} are the linear-relaxation cuts, and integer
 reasoning beyond them for contradictions of parity type; choose orders that
 maximize cancellation; test formulas with planted aligned groups, where
 cancellation is partial.
 \item \textbf{A dichotomy at each node.} Make the three kinds of nodes of
 Section~\ref{sec:approach} precise, so that every node offers either a strong
 clause, a decisive sum, or an aligned structure for pair decisions. By
 Section~\ref{sec:barrier}, a natural dividing line is whether the residual
 formula has few or many strict local maxima, with one-sided pruning as a
 second obstacle.
 \item \textbf{Pair decisions.} By Proposition~\ref{prop:barrier} they are the
 only rule here that can cut strict local maxima. Find orders in which the
 class conditions of Proposition~\ref{prop:chain} are checked early; decide
 whether class decisions can reach the bound $\Ostar((n+1)^t)$ of exact
 interaction classes~\cite{N5} and degrade gracefully when the classes are only
 approximate; and check whether Hastings' sparse formulas with many maxima under
 moves of several variables~\cite{Has16} also defeat class flips.
 \item \textbf{Strategies.} Test the coverage strategy of
 Section~\ref{sec:strategies}, and the largest-ratio choice of
 Section~\ref{sec:availability} on formulas with partial gaps.
\end{enumerate}

\begin{thebibliography}{9}
\raggedright
\bibitem{N1} A literal-deficit measure for Max 2-SAT. Working note,
 22 September 2026 (\texttt{max2sat\_literal\_deficit\_proof.tex}).
\bibitem{N2} Exact Max 2-SAT with an absolute linear polarity gap. Working
 note, 23 September 2026 (\texttt{max2sat\_absolute\_linear\_gap\_proof.tex}).
\bibitem{N3} Max 2-SAT and the domination rule: results and ideas so far.
 Working note, 5 October 2026 (\texttt{max2sat\_results\_and\_ideas.tex}).
\bibitem{N5} Exploring the branching frontier for Max 2-SAT: full clause
 phases, joint completion bounds, and coverage. Draft, 7 October 2026
 (\texttt{max2sat\_branching\_frontier.tex}).
\bibitem{Erd45} P.~Erd\H{o}s. On a lemma of Littlewood and Offord.
 \emph{Bulletin of the American Mathematical Society} 51(12):898--902, 1945.
\bibitem{GK14} A.~Golovnev and K.~Kutzkov. New exact algorithms for the
 2-constraint satisfaction problem. \emph{Theoretical Computer Science}
 526:18--27, 2014.
\bibitem{Has16} M.~B.~Hastings. Local maxima and improved exact algorithm for
 MAX-2-SAT. arXiv:1610.07100, 2016.
\bibitem{AV94} D.~Aldous and U.~Vazirani. ``Go with the winners''
 algorithms. In \emph{Proc.\ 35th FOCS}, pp.~492--501, 1994.
\bibitem{Knu75} D.~E.~Knuth. Estimating the efficiency of backtrack
 programs. \emph{Mathematics of Computation} 29(129):121--136, 1975.
\end{thebibliography}

\end{document}
